% ===== §3 Definition and Conceptual Demarcation =====
\section{Definition and Conceptual Demarcation}
\label{p18:sec:definition}

\noindent
The descent of \xref{p18:sec:typology} cleared a place beneath the established injustices and named its register ontological. This section fills the place with a concept exact enough to be defended---which means exact enough to be \emph{bordered}, for a new category earns its keep not by the ground it claims but by the lines it holds against the neighbours that would absorb it. A category of injustice that blurred into exploitation, into the freezing of power, into the constitutive injustice it descends from, or into mere maldistribution would be a redundancy however vividly named. The section therefore states the definition clause by clause, then walks the four borders along which confusion threatens, showing at each that the line is principled and not stipulative. At two of these borders the most developed rival theories---the Marxian and Roemerian analyses of exploitation, the capability analyses of Sen and Nussbaum---press hardest, and the section meets them there rather than gesturing past them.

\subsection{The definition}
\label{p18:sub:def-statement}

\begin{claim}[Generative effacement]
\emph{Generative effacement} is the injustice in which, within a field of co-generation, one party's generative activity occupies the shared generative space in such a way that the other loses the possibility of continuing to constitute herself as a generative being---so that what is harmed is not her holdings, her standing, or her knowability, but her being itself, \emph{qua} generative being.
\end{claim}

Each clause is load-bearing, and the definition is read by weighing them in turn. \emph{Within a field of co-generation}: the wrong presupposes the relational ontology in which value is generated between the parties rather than held by either, and is possible only because the generative space is \emph{shared}---were each generating in isolation, one party's productivity could not occupy the other's room, there being no common room to occupy. \emph{One party's generative activity occupies the shared space}: the agent of harm is not a vice but an activity, and a good one; the verb is \emph{occupies}, not \emph{seizes}, because the stronger pole does not remove something the weaker had but fills, by its own generating, the space in which the weaker's generating would have occurred. \emph{The other loses the possibility of continuing to constitute herself}: the injury is the foreclosure of an ongoing activity, marked by \emph{continuing}---the subject is already constituted and what is lost is the continuation in which her being consists. \emph{What is harmed is her being itself, qua generative being}: the qualifier is exact; the claim is not that she ceases to exist in every sense---she persists biologically, socially, legally---but that she ceases, within the relation, to exist as a generative being, the mode of existence the series holds to be at stake in intimate life.

\subsection{First border: against exploitation}
\label{p18:sub:def-exploitation}

The most tempting reduction, and the one with the most powerful theory behind it, is to exploitation. The temptation is strong because from outside the two look identical: in both, one party flourishes while another does not. The reduction nonetheless fails, and tracing why against the major accounts of exploitation shows the failure to be structural rather than a matter of emphasis.

On the classical Marxian account, exploitation is the appropriation of the surplus the worker produces above what is returned to her; its mechanism is a \emph{division} of her product into a part she retains and a part the other takes. Roemer, dispensing with the labour theory of value, recasts exploitation in terms of property relations: a coalition is exploited if it would be better off withdrawing with its per-capita share of society's alienable assets \parencite{roemer1982}; and Vrousalis, pressing past Roemer's distributive reduction, grounds the complaint in a Non-Servitude Proviso violated by the power-induced extraction of another's labour \parencite{vrousalis2021}. These accounts differ sharply over what exploitation \emph{is}---surplus-extraction, unequal asset-withdrawal, servitude---but they share, without exception, a structural presupposition that effacement violates: \emph{exploitation requires the exploited party to keep producing}. It is her continued production, appropriated or extracted or coerced, that the exploiter lives upon; an exploitation that stopped its victim's production would abolish its own source. The serf must work the lord's demesne; the worker must labour the surplus hours; the servant must serve. Exploitation is parasitic upon an ongoing generativity it does not extinguish but feeds on.

Effacement does the reverse. It does not appropriate the other's production; it \emph{forecloses} it. The effaced party is not a producer whose output is siphoned but a would-be producer whose production never occurs, because the space in which it would have occurred is already filled. Where exploitation says, in effect, \emph{keep generating, that I may take it}, effacement says \emph{do not generate, for I have generated in your place}. The contrast is exact and it is total: the exploited subject is kept generating and robbed; the effaced subject is not robbed but stilled. Vrousalis's proviso sharpens the point rather than dissolving it---effacement is not servitude, for the effaced party is set to no task; she is set to no task at all, which is precisely the harm.

\begin{claim}[Exploitation feeds on generation; effacement forecloses it]
Exploitation and effacement are not variants but structural opposites. Every developed account of exploitation---Marxian surplus-extraction, Roemerian asset-withdrawal, Vrousalis's servitude---presupposes that the wronged party \emph{continues to generate}, since it is her continued generation that is appropriated, extracted, or coerced. Effacement presupposes the contrary: that the wronged party's generation has been \emph{foreclosed}, the space of it pre-filled. One cannot exploit a subject one has effaced, for there is nothing left to appropriate; one cannot efface a subject one exploits, for one needs her producing. The border is the difference between living on another's generativity and extinguishing it.
\end{claim}

\subsection{Second border: against the freezing of power}
\label{p18:sub:def-freezing}

The second border separates effacement from the solidification of an asymmetry into a fixed, unrebalanceable order---the freezing of power. The freezing presupposes a dynamic that is then arrested: two parties exercise power upon one another, and the wrong is that the play between them ossifies into permanent settlement. Effacement is prior to this. It does not freeze an existing dynamic between two generating poles; it forecloses the second pole's generation before any dynamic between equals can form. The frozen relation still has two agents, locked against each other; the effaced relation has, generatively, one, the other having been pre-empted out of the play entirely. The freezing of power is a pathology of an ongoing contest; effacement is the absence of the contestant. A relation can be both---a frozen power-order can also efface---but the concepts come apart: one can freeze without effacing (two agents permanently locked, both still generating within their fixed positions) and efface without freezing (a fluid, loving, unfixed relation in which one pole simply never generates).

\subsection{Third border: against constitutive injustice}
\label{p18:sub:def-constitutive}

The third border is the most delicate, because effacement descends from constitutive injustice and shares its ontological reach. The line, drawn in \xref{p18:sub:typ-descent}, is restated here as a border: constitutive injustice wrongs the subject at the \emph{threshold of formation}, denying the not-yet-subject her arrival; effacement wrongs the \emph{already-formed} subject, extinguishing the activity in which her achieved being consists. They are kin---both ontological, both deeper than distribution, recognition, or epistemic standing---but they fall on opposite sides of the moment of formation. The constitutive wrong concerns whether a subject will come to be; the effacing wrong, whether an existing subject will continue to be. One bars the door to one who has not entered; the other empties the room of one already inside. The kinship is what makes the border worth drawing: without it, effacement collapses upward into its parent and forfeits its distinctness as a wrong done to maturity rather than to formation.

\subsection{Fourth border: against distributive injustice and the capability turn}
\label{p18:sub:def-distributive}

The fourth border is the plainest to state and the most insistently tempting in practice, because effacement so naturally invites description as the effaced party's having received \emph{too little}. But ``too little'' is distributive language, and it misdescribes the harm---and here the capability theorists must be met, for they offer the subtlest version of the distributive reduction and the one that comes nearest the truth.

Sen and Nussbaum, dissatisfied with distributions of resources, relocate justice's concern to \emph{capabilities}: what a person is actually able to do and to be, the real freedoms she enjoys to achieve functionings she has reason to value \parencite{sen1999, nussbaum2000}. This is a genuine advance toward the harm, for it shifts attention from holdings to doings, and effacement is a harm of foreclosed \emph{doing}. One might think the capability approach already names it: is the effaced party not simply a person whose capability to generate has been removed? The answer marks the border precisely. The capability approach treats capabilities as goods to be \emph{secured} for a person---distributed, protected, enabled by institutions---and so remains, for all its advance, within the paradigm of provision: justice ensures she \emph{has} the capability. But effacement is not the absence of a capability the effaced party lacks and should be given; she retains every capability, in full. What is foreclosed is not her capability to generate but her \emph{generating}---the exercise, not the capacity. And the exercise is not a good that can be secured for her, because it is not a good she can be \emph{given} at all: to hand her generation, ready-made, is precisely to confirm its foreclosure, since what arrives is the stronger pole's generation and not hers. The capability approach can secure that she is \emph{able} to generate; it has no resources to ensure that she \emph{does}, for the doing is hers alone and cannot be provisioned. This is the exact sense in which the harm is of being and not of having: distribution and capability alike concern what can be allotted to a subject, and the effaced subject's loss is of an activity that, by its nature, no allotment reaches.

\begin{claim}[Being, not having; exercise, not capacity]
Effacement is irreducible to maldistribution, including its most refined form, the deprivation of capability. The distributive paradigm, even when it advances from resources to capabilities, concerns what can be \emph{secured for} a subject---holdings, functionings, real freedoms. Effacement forecloses not a capability the subject lacks but the \emph{exercise} of capabilities she fully retains; and an exercise cannot be secured for her, since to supply her generation ready-made is to confirm the foreclosure rather than repair it. The border between distribution and effacement is the border between having and being, and, within doing, between the capacity that can be provisioned and the exercise that cannot.
\end{claim}

\subsection{The fixed point for what follows}
\label{p18:sub:def-shape}

With the concept defined and its four borders held, the paper's remaining work is fixed, and the definition is the fixed point against which each later part is measured. \xref{p18:sec:criterion} must give a criterion that detects \emph{this} harm---the foreclosure of the weaker pole's own generative continuation---in the very cases where the four shallower diagnoses report nothing; a criterion detecting anything else would detect the wrong thing. \xref{p18:sec:reticence} must derive the justice that repairs \emph{this} harm; a justice repairing anything else would repair the wrong thing. And \xref{p18:sec:ethics} must ground an obligation to cease \emph{this} occupation; an ethics grounding anything else would ground the wrong obligation. Everything that follows is the working-out of the harm bordered here.
